\chapter{How Do I Pay for Healthcare}

\section*{Definition and Overview}
In Australia, healthcare is paid for through a mix of public and private arrangements. Medicare, the national public health scheme, covers free treatment as a public patient in public hospitals and subsidises GP visits and many medical services. The Pharmaceutical Benefits Scheme (PBS) subsidises prescription medicines, and private health insurance covers private hospital treatment and some out-of-hospital services. Understanding these arrangements helps people access care, manage costs, and avoid unexpected bills. This chapter explains how the system works, what is free, what costs money, and where to get help with costs.

\section*{Epidemiology}
Most Australians are enrolled in Medicare, which is funded through the Medicare levy on income. Around half of Australians hold private health insurance, most commonly for private hospital cover and extras such as dental and physiotherapy. Each year, Australians pay billions in out-of-pocket costs for GP visits, dental care, and medicines, and cost is a recognised barrier to care for some people. Public hospital emergency care is free, and people in financial hardship can access concessions for medicines and health services.

\section*{Aetiology and Pathophysiology}
\begin{itemize}
    \item \textbf{Medicare:} covers public hospital treatment, subsidised GP and specialist visits, and some tests; funded by the Medicare levy
    \item \textbf{Bulk billing:} when a doctor bills Medicare directly, the patient pays nothing
    \item \textbf{Out-of-pocket costs:} when providers charge more than the Medicare rebate, the patient pays the gap
    \item \textbf{PBS:} subsidises prescription medicines, with a general co-payment and lower concessional rates
    \item \textbf{Private health insurance:} covers private hospital costs and some extras, with premiums, excesses, and exclusions
    \item \textbf{Public hospitals:} free as a public patient, with waiting lists; private patients pay via insurance or self-funding
    \item \textbf{Concessions:} Health Care Cards and other schemes reduce costs for people on low incomes
\end{itemize}

\section*{Clinical Features}
\begin{itemize}
    \item GP visits: covered by Medicare, often bulk billed, sometimes with a gap
    \item Specialist visits: partially covered by Medicare, usually with a gap
    \item Hospital care: free as a public patient; private care funded by insurance
    \item Medicines: PBS-subsidised prices at pharmacies
    \item Dental, optical, and allied health: generally not covered by Medicare, with some insurance extras
    \item Emergency care: free in public hospital emergency departments
    \item Ambulance services: covered in some states and by some insurers, with costs elsewhere
\end{itemize}

\section*{Differential Diagnosis}
\begin{itemize}
    \item Costs differ between bulk-billed and gap-paying providers, so shopping around helps
    \item Private insurance choices vary in premiums, excess, and coverage
    \item Medicare Safety Nets reduce costs once annual thresholds are reached
    \item State and territory differences affect ambulance and some services
    \item Concession cards reduce PBS and other costs for eligible people
\end{itemize}

\section*{When to Seek Medical Care}
Everyone should know how to access care: GPs for routine care, emergency departments for urgent problems, and healthdirect for advice. Never delay needed care because of cost concerns; services are available to help.

\textbf{Call 000 (or local emergency services) for:}
\begin{itemize}
    \item Chest pain, severe breathlessness, or signs of stroke
    \item Severe bleeding, major injury, or collapse
    \item Any life-threatening emergency; care is not withheld based on ability to pay
    \item Severe allergic reactions and seizures
\end{itemize}

\section*{Diagnosis and Investigations}
\begin{itemize}
    \item \textbf{Checking Medicare:} confirm enrolment and eligibility, including through the Medicare app
    \item \textbf{Checking insurance:} understand what your private policy covers before needing it
    \item \textbf{Asking providers:} always ask about costs before treatment and whether bulk billing is available
    \item \textbf{Understanding thresholds:} PBS Safety Net and Medicare Safety Net reduce annual costs
    \item \textbf{Seeking concessions:} apply for Health Care Cards and other concessions if eligible
\end{itemize}

\section*{Management}
\begin{itemize}
    \item \textbf{Use bulk-billed services:} many GPs, clinics, and pathology services bulk bill
    \item \textbf{Public hospital care:} free treatment as a public patient
    \item \textbf{PBS medicines:} prescribed medicines are subsidised; ask about cheaper options
    \item \textbf{Private insurance:} choose cover that fits your needs and use it for extras and private care
    \item \textbf{Community health centres:} low-cost care in many areas
    \item \textbf{Financial counselling:} hospitals and community services help with costs and hardship
    \item \textbf{Prevention:} free and subsidised screening and immunisation reduce future costs
\end{itemize}

\section*{Complications}
\begin{itemize}
    \item Delayed care because of cost concerns, leading to worse health outcomes
    \item Unexpected out-of-pocket bills, especially from private hospitals and specialists
    \item Unaffordable medicines leading to non-adherence
    \item Gaps in dental and allied health care
    \item Financial stress and debt from healthcare costs
    \item Confusion about entitlements, with people missing benefits they are owed
\end{itemize}

\section*{Prognosis and Outcomes}
Australia's mixed public-private system provides broad access to care for most people, and out-of-pocket costs are modest for most services. People who understand their entitlements, ask about costs, and use bulk billing and concessions pay less and access care earlier. Financial assistance is available, and no one is denied emergency care. With informed use, the system delivers high-quality care at a cost that is manageable for the great majority.

\section*{Prevention and Self-Care}
\begin{itemize}
    \item Enrol in Medicare and understand your cover
    \item Ask every provider about costs and bulk billing before treatment
    \item Compare private insurance policies before buying
    \item Use Medicare Safety Net and PBS Safety Net thresholds
    \item Apply for concession cards if eligible
    \item Use free services: public hospitals, community health, and healthdirect
    \item Attend free screening and immunisation programs
    \item Seek financial counselling if healthcare costs cause hardship
\end{itemize}

\section*{Sources and Further Reading}
The following reputable sources provide further information for healthcare students and patients seeking reliable, current advice about paying for healthcare.

\begin{itemize}
    \item Services Australia. \textit{Medicare and your health care costs.} https://www.servicesaustralia.gov.au/
    \item Healthdirect. \textit{Costs of health care in Australia.} https://www.healthdirect.gov.au/
    \item Private Health Insurance Ombudsman. \textit{Comparing and understanding private health insurance.} https://www.privatehealth.gov.au/
    \item Australian Government Department of Health. \textit{Pharmaceutical Benefits Scheme information.} https://www.health.gov.au/
    \item Medicare Australia. \textit{Medicare benefits and safety nets.} https://www.servicesaustralia.gov.au/
    \item National Health Funding Pool. \textit{How public hospitals are funded.} https://www.publichospitalfunding.gov.au/
\end{itemize}
